% GENERATED FILE. Edit canonical lessons, then run npm run generate:books.
% canonical-source-hash: fnv1a64:5a188aef66849c57
% canonical-lessons: RU-C207-bolet, RU-C207-lechit, RU-C207-dyshat, RU-C207-kashlyat, RU-C207-chikhat

\chapter{To be ill, To treat, To breathe, To cough, To sneeze}
\label{ch:ru-to-be-ill-to-treat-to-breathe-to-cough-to-sneeze}
\hlchaptermodality{207}

\begin{chapteropening}
\textbf{By the end of this chapter:} \emph{I can say to be ill, to treat, to breathe, to cough and to sneeze.}

Complete the last lesson of chapter 207: I can say to be ill, to treat, to breathe, to cough and to sneeze.
\end{chapteropening}

\section[\texorpdfstring{\ru{болеть} (bolét)}{болеть (bolét)}]{\ru{болеть} (bolét) — to be ill; to ache}

\label{lesson:RU-C207-bolet}



\begin{quote}
Before the new one: say the Russian for to paint, then the Russian for to clean.
\end{quote}

\subsection*{You'll want to know: \ru{болеть}}
\textbf{\ru{болеть}} (\emph{bolét}) — \textquotedblleft{}to be ill; to ache\textquotedblright{}.

\begin{grammarlens}[title={What you've built}]
One more word for this chapter.
\end{grammarlens}

\subsection*{Your turn}
\begin{itemize}
\raggedright
  \item \emph{Say it:} \emph{bolét}
  \item \emph{Say it:} \emph{bolét}, once more
  \item \emph{Say it:} say \emph{bolét}
  \item \emph{Recall:} say the Russian for to paint, then the Russian for to clean, then say \emph{bolét} again
\end{itemize}

\begin{tcolorbox}[breakable,colback=teal!4,colframe=teal!35!black,title={Before you move on}]
What does \ru{болеть} mean? (\textquotedblleft{}To be ill; to ache\textquotedblright{}.) Say it once more.
\end{tcolorbox}
\section[\texorpdfstring{\ru{лечить} (lechít)}{лечить (lechít)}]{\ru{лечить} (lechít) — to treat (an illness)}

\label{lesson:RU-C207-lechit}



\begin{quote}
Before the new one: say the Russian for to clean, then the Russian for to be ill; to ache.
\end{quote}

\subsection*{You'll want to know: \ru{лечить}}
\textbf{\ru{лечить}} (\emph{lechít}) — \textquotedblleft{}to treat (an illness)\textquotedblright{}.

\begin{grammarlens}[title={What you've built}]
One more word for this chapter.
\end{grammarlens}

\subsection*{Your turn}
\begin{itemize}
\raggedright
  \item \emph{Say it:} \emph{lechít}
  \item \emph{Say it:} \emph{lechít}, once more
  \item \emph{Say it:} say \emph{lechít}
  \item \emph{Recall:} say the Russian for to clean, then the Russian for to be ill; to ache, then say \emph{lechít} again
\end{itemize}

\begin{tcolorbox}[breakable,colback=teal!4,colframe=teal!35!black,title={Before you move on}]
What does \ru{лечить} mean? (\textquotedblleft{}To treat (an illness)\textquotedblright{}.) Say it once more.
\end{tcolorbox}
\section[\texorpdfstring{\ru{дышать} (dyshát)}{дышать (dyshát)}]{\ru{дышать} (dyshát) — to breathe}

\label{lesson:RU-C207-dyshat}



\begin{quote}
Before the new one: say the Russian for to be ill; to ache, then the Russian for to treat.
\end{quote}

\subsection*{You'll want to know: \ru{дышать}}
\textbf{\ru{дышать}} (\emph{dyshát}) — \textquotedblleft{}to breathe\textquotedblright{}.

\begin{grammarlens}[title={What you've built}]
One more word for this chapter.
\end{grammarlens}

\subsection*{Your turn}
\begin{itemize}
\raggedright
  \item \emph{Say it:} \emph{dyshát}
  \item \emph{Say it:} \emph{dyshát}, once more
  \item \emph{Say it:} say \emph{dyshát}
  \item \emph{Recall:} say the Russian for to be ill; to ache, then the Russian for to treat, then say \emph{dyshát} again
\end{itemize}

\begin{tcolorbox}[breakable,colback=teal!4,colframe=teal!35!black,title={Before you move on}]
What does \ru{дышать} mean? (\textquotedblleft{}To breathe\textquotedblright{}.) Say it once more.
\end{tcolorbox}
\section[\texorpdfstring{\ru{кашлять} (káshlyat)}{кашлять (káshlyat)}]{\ru{кашлять} (káshlyat) — to cough}

\label{lesson:RU-C207-kashlyat}



\begin{quote}
Before the new one: say the Russian for to treat, then the Russian for to breathe.
\end{quote}

\subsection*{You'll want to know: \ru{кашлять}}
\textbf{\ru{кашлять}} (\emph{káshlyat}) — \textquotedblleft{}to cough\textquotedblright{}.

\begin{grammarlens}[title={What you've built}]
One more word for this chapter.
\end{grammarlens}

\subsection*{Your turn}
\begin{itemize}
\raggedright
  \item \emph{Say it:} \emph{káshlyat}
  \item \emph{Say it:} \emph{káshlyat}, once more
  \item \emph{Say it:} say \emph{káshlyat}
  \item \emph{Recall:} say the Russian for to treat, then the Russian for to breathe, then say \emph{káshlyat} again
\end{itemize}

\begin{tcolorbox}[breakable,colback=teal!4,colframe=teal!35!black,title={Before you move on}]
What does \ru{кашлять} mean? (\textquotedblleft{}To cough\textquotedblright{}.) Say it once more.
\end{tcolorbox}
\section[\texorpdfstring{\ru{чихать} (chikhát)}{чихать (chikhát)}]{\ru{чихать} (chikhát) — to sneeze}

\label{lesson:RU-C207-chikhat}



\begin{quote}
Before the new one: say the Russian for to breathe, then the Russian for to cough.
\end{quote}

\subsection*{You'll want to know: \ru{чихать}}
\textbf{\ru{чихать}} (\emph{chikhát}) — \textquotedblleft{}to sneeze\textquotedblright{}.

\begin{grammarlens}[title={What you've built}]
One more word for this chapter.
\end{grammarlens}

\subsection*{Your turn}
\begin{itemize}
\raggedright
  \item \emph{Say it:} \emph{chikhát}
  \item \emph{Say it:} \emph{chikhát}, once more
  \item \emph{Say it:} say \emph{chikhát}
  \item \emph{Recall:} say the Russian for to breathe, then the Russian for to cough, then say \emph{chikhát} again
\end{itemize}

\begin{tcolorbox}[breakable,colback=teal!4,colframe=teal!35!black,title={Before you move on}]
What does \ru{чихать} mean? (\textquotedblleft{}To sneeze\textquotedblright{}.) Say it once more.
\end{tcolorbox}
